\clearpage
\fancyhead[L]{\small JUNO.}
\hypertarget{native-261-lower}{}
\begin{center}
JUNO.

\rule{0.75\linewidth}{0.4pt}

Monatliche Correspondenz. Vol. X. pp. 463–471. November 1804.

\rule{0.75\linewidth}{0.4pt}
\end{center}

\noindent—\quad—\quad—\par
Dr. \textsc{Gauss} asked us to determine the position of a star that was very close to \emph{Juno} on 28 Sept. and with which he had compared the planet. By averaging observations made over five days, we found the following position for this star. It may also be useful to other observers.

\begin{center}
\begingroup\fontsize{9}{11}\selectfont\setlength{\tabcolsep}{3pt}\renewcommand{\arraystretch}{1.1}
\begin{tabular}{r|r|r|r|r}
Magnitude & \shortstack{Mean right\\ascension\\1804} & \shortstack{Annual\\change} & \shortstack{Mean southern\\declination\\1804} & \shortstack{Annual\\change}\\ \hline
8 & $358^\circ\ 19'\ 34''42$ & $+46''1$ & $5^\circ\ 13'\ 57''27$ & $-20''0$ \\
\end{tabular}\endgroup
\end{center}


Dr. \textsc{Gauss} continues tirelessly with the difficult and laborious work of improving the elements of this new planet’s orbit. “Since my last letter,” (writes this master of calculation) “I have continued to observe our \emph{Juno} diligently, but I no longer use my own observations for the elements, since those you were kind enough to send me are already beginning to cover a considerable arc. Meanwhile I tried to refine Juno’s orbit again by taking the mean motion of Ceres and Pallas as a hypothetical basis for Juno’s motion too. The attempt failed. I then made a second determination of the elements independently of hypotheses. These are the results.”

\begin{center}
\emph{II\textsuperscript{nd} Elements of Juno.}
\end{center}

\begin{center}
\begingroup\fontsize{9}{11}\selectfont\setlength{\tabcolsep}{3pt}\renewcommand{\arraystretch}{1.1}
\begin{tabular}{lr}
Epoch 1804, 30 Sept., at $0^{\mathrm u}$ mean time  &  \\
\quad in Seeberg . . . . . . . . . & $21^\circ\ 17'\ 47''$ \\
Daily mean motion . . . . & $836''89$ \\
Aphelion . . . . . . . . . & $231^\circ\ 38'\ 1''$ \\
Eccentricity . . . . . . . . . & $0.254964$ \\
Logarithm of the semiaxis . . . . . & $0.4182255$ \\
Node . . . . . . . . . . . & $170^\circ\ 46'\ 41''$ \\
Inclination of the orbit . . . . . . . & $12^\circ\ 19'\ 43''$ \\
\end{tabular}\endgroup
\end{center}


“If no errors have crept into my calculations, it can already be taken as fairly settled that Juno’s mean motion is \emph{not equal} to that of Ceres and Pallas. It is \emph{considerably greater}, so its period of revolution and mean distance from the Sun are smaller. Continued observations will very soon give us certainty on this point.”

At first it seemed as though Dr. \textsc{Gauss} would find the same period of revolution for \emph{Juno} as for Ceres and Pallas. This would have been quite at odds with Dr. \textsc{Olbers}’s hypothesis about the origin of this small plan%
\SourceLeaf{262}%
et. For, on the assumption that all these small planets might be only debris and fragments of a larger planet that was destroyed, it can indeed be shown that the periods of revolution of these small fragments cannot be \emph{very different}. But it was highly unlikely that all the pieces of the shattered planet should have acquired exactly the same speed. Now, as we have seen above, our excellent \textsc{Gauss} finds the mean motion of \emph{Juno} much faster than that of Ceres and Pallas.

Dr. \textsc{Olbers} puts it as follows. “The overall position of Juno’s orbit is consistent with my hypothesis (which, incidentally, I still intend to put forward only as a hypothesis). Its nodes with Ceres’s orbit now lie about 24 degrees from the node of Pallas’s orbit. But with the inclinations of these orbits already so different, Jupiter’s attraction must cause the nodes to shift unevenly. Juno’s orbit now lies well inside Ceres’s orbit at the descending node on Ceres’s orbit, where Pallas’s orbit comes so close to it. The aphelia of all these orbits move quite differently from the nodes, so the positions of the lines of apsides relative to the nodes keep changing. And the orbits have nearly equal axes and very unequal eccentricities. It follows that these orbits will actually intersect at certain times and have actually intersected in the past. If, for example, I adopt the annual shifts of the aphelia determined by \textsc{Oriani}, $106''1$ for Pallas and $120''9$ for Ceres, and take the nodes to be fixed relative to the stars and the inclinations to be invariable, it follows that the orbits of Ceres and Pallas actually intersected 7463 years ago at the descending node of Pallas on Ceres’s orbit and will intersect again in 282 years. An intersection of this kind will occur at the ascending node in 925 years. Just as Pallas’s orbit now lies inside Ceres’s at both nodes, in 1000 years Ceres’s orbit will lie inside Pallas’s. These considerations, however, can yield nothing decisive until the perturbations of all three orbits have been fully worked out.”

\Needspace{24\baselineskip}
Until that can be done, Dr. \textsc{Gauss} is continuing to improve the elliptic elements of the orbit and compare them with the best meridian observations. Here is such a comparison of all these observations with his second elements of \emph{Juno}.

\begin{center}
\begingroup\fontsize{8}{10}\selectfont\setlength{\tabcolsep}{3pt}\renewcommand{\arraystretch}{1.1}
\begin{tabular}{r|r|r|r|r|l}
 & \multicolumn{2}{c|}{Calculated} & \multicolumn{2}{c|}{Difference} & \\ 1804 & \shortstack{Right\\ascension} & \shortstack{Southern\\declination} & \shortstack{in R.A.} & \shortstack{in decl.} & Observer \\ \hline
13 Sept. & $0^\circ\ 44'\ 56''0$ & $1^\circ\ 52'\ 37''6$ & $-0''4$ & $+0''6$ & von Zach \\
14 & $0\ 35\ 42.9$ & $2\ 5\ 40.3$ & $+1.9$ & $+4.8$ & — \\
15 & $0\ 26\ 19.8$ & $2\ 18\ 48.6$ & $-0.2$ & $+10.1$ & — \\
17 & $0\ 7\ 8.0$ & $2\ 45\ 18.8$ & $-0.1$ & $-5.7$ & — \\
18 & $359\ 57\ 20.9$ & $2\ 58\ 38.7$ & $+0.4$ & $+7.8$ & — \\
20 & $359\ 37\ 39.0$ & $3\ 25\ 24.7$ & $.\ .\ .\ .$ & $+21.2$ & — \\
23 & $359\ 7\ 17.4$ & $4\ 5\ 36.5$ & $+2.5$ & $+2.0$ & — \\
23 & $359\ 7\ 3.3$ & $4\ 5\ 55.0$ & $+3.3$ & $+8.0$ & Dr. Burckhardt \\
25 & $358\ 46\ 45.8$ & $4\ 32\ 41.4$ & $-2.2$ & $+13.2$ & Dr. Maskelyne \\
27 & $358\ 26\ 56.9$ & $4\ 58\ 43.7$ & $-3.7$ & $-12.8$ & von Zach \\
28 & $358\ 16\ 58.7$ & $5\ 11\ 49.3$ & $-7.9$ & $+5.5$ & — \\
29 & $358\ 6\ 48.0$ & $5\ 25\ 11.7$ & $-4.3$ & $+15.0$ & Dr. Maskelyne \\
30 & $357\ 57\ 18.2$ & $5\ 37\ 50.4$ & $-8.5$ & $+6.7$ & von Zach \\
2 Oct. & $357\ 38\ 6.4$ & $6\ 2\ 59.7$ & $-9.6$ & $+0.2$ & — \\
4 & $357\ 19\ 34.8$ & $6\ 27\ 40.2$ & $-4.7$ & $+1.5$ & — \\
5 & $357\ 10\ 37.3$ & $6\ 39\ 43.7$ & $-4.2$ & $+3.8$ & — \\
6 & $357\ 1\ 52.8$ & $6\ 51\ 35.6$ & $-3.5$ & $+3.3$ & — \\
\end{tabular}\endgroup
\end{center}

\SourceLeaf{263}
\indent Using the same elements, Dr. \textsc{Gauss} has also calculated the following ephemeris for the future course of \emph{Juno} at midnight in Seeberg. [From Oct. 15 to Dec. 2 it gives right ascension and declination.]

Around 18 December, Juno and Ceres come close together as seen from Earth.

\noindent—\quad—\quad—\par

\hfill[\textsc{von Zach.}]

